\chapter{Streaming the multipoint lock-in}
\label{ch:mpstream}

Everything so far has been the two-point method: one pair of parked frequencies on
one flank pair, giving the field projection along a single NV axis. The multipoint
notebook (\file{Modules/Lockin_module.ipynb}) parks 16 frequencies instead, which
is what makes a full three-axis reconstruction possible --- and until 2026-08-17 it
had only the slowest of the three acquisition paths.

This chapter is what changes when \co{stream()} is pointed at it. Nothing in the
FPGA program needs to change: \co{emission_order()}, \co{_fold_slots()} and
\co{stream()} are all written in terms of $n_\text{slots}$, so 16 parked points
work exactly as 2 do. What changes is the arithmetic, and one constraint that has
never bound on the two-point path.

\section{Why it is worth doing}

\co{acquire()} costs \SI{8.5}{}--\SI{11.4}{\milli\second} of host round-trip
whatever the program does (Ch.~\ref{ch:timing}), and Steps 4/4B/4BN call it once
per sample. With 16 points and a reference the FPGA work is
\SI{3.84}{\milli\second} --- comfortably \emph{under} the floor --- so the live
loop is host-bound at

\[
\frac{1}{\max(\SI{10}{\milli\second},\ \SI{3.84}{\milli\second})}\approx\SI{26}{\hertz},
\]

\noindent and no FPGA setting moves it. \co{stream()} pays the host cost once for
the whole run, so the rate becomes the cadence itself:

\begin{center}
\small
\begin{tabular}{@{}lccrrl@{}}
\toprule
\hd{Path} & \hd{slots} & \hd{reads/slot} & \hd{$t_\text{rep}$} & \hd{Rate} & \hd{Note} \\
\midrule
Step 4 / 4B / 4BN (live loop) & 16 & 2 & \SI{3.84}{\milli\second} & \SI{26}{\hertz} & host-bound, not FPGA-bound \\
\textbf{Step 4S}, reference kept & 16 & 2 & \SI{3.84}{\milli\second} & \textbf{\SI{260}{\hertz}} & the default; reconstructable \\
Step 4S, \co{skip_reference}     & 16 & 1 & \SI{1.92}{\milli\second} & \SI{521}{\hertz} & projections only --- see below \\
\bottomrule
\end{tabular}
\end{center}

A factor of ten, and it comes with the two structural properties from
Ch.~\ref{sec:stream}: nothing is re-armed mid-run, so the replay race that costs
burst mode half its batches cannot occur; and the samples are consecutive reps of
one program, so the time axis is the cadence exactly.

\begin{keybox}[The reference readout stays on, and that is not a rate compromise]
It is tempting to set \co{multipoint_skip_reference = True} and double the rate,
as the two-point path does. Do not, for any run you intend to reconstruct.

The vector inversion consumes \co{peak_NN / peak_NN_ref} --- each parked point
normalised by its own off-resonance reference. That division is what puts the
values on the calibration's scale, because the calibration itself is built from
\co{mw_on / median(mw_on)} inside \co{_load_operator_full_scan}. Feed it
unnormalised counts and it returns field values around \num{1e5}\,\si{\micro\tesla},
which is not a subtle failure but is an easy one to mistake for a geometry
problem.

The two-point path can skip the reference because it normalises against a stored
constant from the calibration JSON instead (\S\ref{sec:asm}). The multipoint path
has no equivalent.
\end{keybox}

\section{The constraint that actually binds: buffer headroom}
\label{sec:mpheadroom}

The board's accumulated buffer holds \co{avg_maxlen} entries, one per readout per
rep, and the board worker raises if unread entries ever reach that. The budget is
shared across \emph{every} readout in a shot, so what matters is

\[
\co{shots_that_fit} = \left\lfloor \frac{\co{avg_maxlen}}{n_\text{slots}\times n_\text{readouts per slot}} \right\rfloor .
\]

\begin{center}
\small
\begin{tabular}{@{}lrrl@{}}
\toprule
\hd{Configuration} & \hd{reads/shot} & \hd{relative headroom} & \\
\midrule
two-point, skip reference \emph{(all data so far)} & 2  & $1\times$      & has never come close \\
two-point, with reference                          & 4  & $1/2\times$    & \\
16-point, skip reference                           & 16 & $1/8\times$    & \\
\textbf{16-point, with reference (Step 4S)}        & 32 & $1/16\times$   & \textbf{check before every long run} \\
\bottomrule
\end{tabular}
\end{center}

Sixteen times less room, and the stride the worker transfers in is 10\% of it, so
the margin between ``fine'' and ``the run dies at 20 seconds'' is much narrower
than anything the two-point path has explored. Step 4S calls
\co{stream_headroom()} before starting and warns if fewer than four strides fit.

\begin{openbox}[The headroom number has never actually been read on this rig]
\co{stream_headroom()} was broken for the whole 2026-08-17 session --- it
subscripted the Pyro proxy instead of \co{qd.soccfg} (\S\ref{sec:stride}) --- so
every run printed ``could not read avg_maxlen'' and proceeded blind. It is fixed,
but that means the first Step 4S run is also the first time this rig's
\co{avg_maxlen} will be known. If it turns out to be tight, the levers in order of
preference are: shorten $\tau$ (the buffer is counted in \emph{entries}, so this
does not help --- it only makes the entries arrive faster, which is worse), reduce
the parked-point count, or drop the reference and give up reconstruction.
\end{openbox}

\section{Why the reconstruction is not inline}
\label{sec:mprecon}

Step 4BN reconstructs the field vector inside its acquisition loop, once per
batch. Step 4S does not, and the difference is not stylistic.

At \SI{26}{\hertz} there is \SI{38}{\milli\second} between samples and the fit
costs a few milliseconds; there is room. At \SI{260}{\hertz} there is
\SI{3.84}{\milli\second} between samples and there is not. More importantly, the
consequence of falling behind is different in kind: the live loop simply runs
slower, because nothing is happening on the board between its \co{acquire()}
calls. A stream has the tProc running continuously, so a consumer that cannot keep
up lets unread shots accumulate in the buffer of \S\ref{sec:mpheadroom} until it
overflows and \emph{the run dies}. Slow post-processing inside a stream loop is
not a performance question, it is a correctness one.

So Step 4S writes files and does nothing else, and Step 5S does the work
afterwards, once, over the whole run --- through
\co{nv_toolkit.tui._compute_live_snapshot}, the same entry point the existing
offline reconstruction cell uses. Same calibration, same inversion, so a streamed
run and a live run are directly comparable.

\co{S5_DECIMATE} is there for the case where even the offline fit is too slow to
be comfortable: it reconstructs every $N$th sample. Note that this is a genuine
decimation, unlike the acquisition-side bin of \S\ref{sec:bin} --- it neither
averages nor anti-aliases. Bin in Step 4S; decimate in Step 5S only to make a
first look fast.

\section{What the two cells produce}

\begin{center}
\small
\begin{tabular}{@{}lp{9.6cm}@{}}
\toprule
\hd{File} & \hd{Contents} \\
\midrule
\co{multipoint_lockin_stream_<stamp>.csv} &
  One row per output sample: \co{packet}, \co{rep_index}, \co{reps_averaged},
  \co{time_s}, \co{timestamp_epoch_s}, then \co{peak_NN}, \co{peak_NN_freq_mhz}
  and \co{peak_NN_ref} for each parked point. \\[3pt]
\co{..._wide.csv} &
  \co{time_s} plus \co{peak_NN} normalised by its reference --- the format the
  inversion consumes. Rebuilt by Step 5S rather than trusted, so a run copied in
  from another machine works. \\[3pt]
\co{..._projection_rows.csv} &
  Per-block $\Delta f$ and new peak frequency. \\[3pt]
\co{..._vector_rows.csv} &
  $\Delta B_x$, $\Delta B_y$, $\Delta B_z$ per sample. \\
\bottomrule
\end{tabular}
\end{center}

The naming matches the live-loop convention (\co{multipoint_lockin_live_*} and its
\co{_wide} / \co{_vector_rows} / \co{_projection_rows} companions), so the existing
post-processing and the session folders under \file{data/results/} stay consistent.

\section{What Step 5S checks, and why each check exists}

\begin{center}
\small
\begin{tabular}{@{}lp{9.4cm}@{}}
\toprule
\hd{Check} & \hd{What a failure means} \\
\midrule
\co{rep_index} gaps & A dropped rep. Every timestamp \emph{after} the first gap is
  wrong, because the axis is built by counting reps rather than by measuring time.
  This is the one failure that silently corrupts the frequency scale of every
  spectrum downstream. \\[3pt]
duplicate \co{rep_index} & Would be the burst replay appearing where it should be
  impossible. Has never been seen in a stream. \\[3pt]
timestamp spacing & Should be identical to four decimal places. Spread here means
  the same thing as a rep gap. \\[3pt]
PL droop & A continuous run has no relaxation gap between batches, so the
  common-mode level walks over the run --- $\sim$25\% over \SI{30}{\second} on the
  two-point path. It is common to every parked point, which is what the reference
  normalisation removes. \\[3pt]
rank-deficient samples & A geometry problem with the parked plan, not a streaming
  one: the same plan fails the same way in Step 4B. \\
\bottomrule
\end{tabular}
\end{center}

\begin{openbox}[Not yet run on hardware]
Everything in this chapter is derived from the timing model and from
\co{stream()}'s behaviour on the two-point path, which is well characterised. The
multipoint streaming cells have been checked for syntax and for name resolution
against the notebook's own globals, and nothing more. The first rig run should
compare a Step 4S run against a Step 4BN run taken back to back: the reconstructed
$\Delta B$ should agree within noise, and Step 4S should report ten times the rate
with \co{reps_missing = 0}.
\end{openbox}
